% ============ §2 Modules (原书 pp.6-14) ============
\section{Modules}\label{sec:2}\index{Modules, section 2}

Let $A$ be a ring and $M$ an $A$-module. Given submodules $N$, $N'$
of $M$, the set $\{a \in A \mid aN' \subset N\}$ is an ideal of $A$, which we write $N:N'$
or $(N:N')_A$. Similarly, if $I \subset A$ is an ideal then $\{x \in M \mid Ix \subset N\}$ is a
submodule of $M$, which we write $N:I$ or $(N:I)_A$. For $a \in A$ we define
$N:a$ similarly. The ideal $0:M$ is called the \emph{annihilator}\index{annihilator} of $M$, and written
$\ann(M)$. We can consider $M$ as a module over $A/\ann(M)$. If $\ann(M) = 0$
we say that $M$ is a \emph{faithful} $A$-module.\index{faithful module} For $x \in M$ we write
$\ann(x) = \{a \in A \mid ax = 0\}$.

If $M$ and $M'$ are $A$-modules, the set of $A$-linear maps from $M$ to $M'$
is written $\Hom_A(M, M')$. This becomes an $A$-module if we define the sum
$f + g$ and the scalar product $af$ by
\[
  (f + g)(x) = f(x) + g(x), \qquad (af)(x) = af(x);
\]
(the fact that $af$ is $A$-linear depends on $A$ being commutative).

To say that $M$ is an $A$-module is to say that $M$ is an Abelian group
under addition, and that a scalar product $ax$ is defined for $a \in A$ and
$x \in M$ such that the following hold:
\[
  (*) \quad a(x + y) = ax + ay, \quad (ab)x = a(bx), \quad (a + b)x = ax + bx, \quad 1x = x;
\]
for fixed $a \in A$ the map $x \mapsto ax$ is an endomorphism of $M$ as an additive
group. Let $E$ be the set of endomorphisms of the additive group $M$;
defining the sum and product of $\lambda, \mu \in E$ by
\[
  (\lambda + \mu)(x) = \lambda(x) + \mu(x), \qquad (\lambda\mu)(x) = \lambda(\mu(x))
\]
makes $E$ into a ring (in general non-commutative), and giving $M$ an
$A$-module structure is the same thing as giving a homomorphism
$A \to E$. Indeed, if we write $a_L$ for the element of $E$ defined by $x \mapsto ax$ then
$(*)$ becomes
\[
  (ab)_L = a_Lb_L, \qquad (a + b)_L = a_L + b_L, \qquad (1)_L = 1_L.
\]
We can express the fact that $\varphi\colon M \to M$ is $A$-linear by saying that
$\varphi \in E$ and that $\varphi$ commutes with $a_L$ for $a \in A$, that is $a_L\varphi = \varphi a_L$. Since
$A$ is commutative, $a_L$ is itself an $A$-linear map of $M$ for $a \in A$. We normally
write simply $a\colon M \to M$ for the map $a_L$.

If $M$ is a $B$-module and $f\colon A \to B$ a ring homomorphism, then
we can make $M$ into an $A$-module by defining $a \cdot x = f(a) \cdot x$ for $a \in A$ and
$x \in M$. This is the $A$-module structure defined by the composite of
$f\colon A \to B$ with $B \to E$, where $E$ is the endomorphism ring of the additive
group of $M$, and $B \to E$ is the map defining the $B$-module structure of $M$.

If $M$ is finitely generated as an $A$-module we say simply that $M$ is a
\emph{finite $A$-module}\index{finite module}, or is \emph{finite over $A$}. A standard technique applicable to
finite $A$-modules is the `determinant trick', one form of which is as follows
(taken from Atiyah and Macdonald~\cite{AM}).

\begin{theorem}\label{thm:2.1}
Suppose that $M$ is an $A$-module generated by $n$ elements,
and that $\varphi \in \Hom_A(M, M)$; let $I$ be an ideal of $A$ such that $\varphi(M) \subset IM$.
Then there is a relation of the form
\[
  (**) \quad \varphi^n + a_1\varphi^{n-1} + \cdots + a_{n-1}\varphi + a_n = 0,
\]
with $a_i \in I^i$ for $1 \leq i \leq n$ (where both sides are considered as endomorphisms
of $M$).
\end{theorem}

\begin{proof}
Let $M = A\omega_1 + \cdots + A\omega_n$; by the assumption $\varphi(M) \subset IM$ there exist
$a_{ij} \in I$ such that $\varphi(\omega_i) = \sum_{j=1}^{n} a_{ij}\omega_j$. This can be rewritten
\[
  \sum_{j=1}^{n} (\varphi\delta_{ij} - a_{ij})\omega_j = 0 \quad (\text{for } 1 \leq i \leq n),
\]
where $\delta_{ij}$ is the Kronecker symbol. The coefficients of this system of linear
equations can be viewed as a square matrix $(\varphi\delta_{ij} - a_{ij})$ of elements of $A'[\varphi]$,
the commutative subring of the endomorphism ring $E$ of $M$ generated by
the image $A'$ of $A$ together with $\varphi$; let $b_{ij}$ denote its $(i,j)$th cofactor, and
$d$ its determinant. By multiplying the above equation through by $b_{ik}$ and
summing over $i$, we get $d\omega_k = 0$ for $1 \leq k \leq n$. Hence $d \cdot M = 0$, so that $d = 0$
as an element of $E$. Expanding the determinant $d$ gives a relation of the
form $(**)$.
\end{proof}

\begin{remark}
As one sees from the proof, the left-hand side of $(**)$ is the
characteristic polynomial of $(a_{ij})$,
\[
  f(X) = \det(X\delta_{ij} - a_{ij})
\]
with $\varphi$ substituted for $X$. If $M$ is the free $A$-module with basis
$\omega_1, \ldots, \omega_n$, and $I = A$, the above result is nothing other than the classical
Cayley--Hamilton theorem: let $f(X)$ be the characteristic polynomial of
the square matrix $\varphi = (a_{ij})$; then $f(\varphi) = 0$.
\end{remark}

\begin{theorem}[NAK]\label{thm:2.2}
Let $M$ be a finite $A$-module and $I$ an ideal of $A$. If
$M = IM$ then there exists $a \in A$ such that $aM = 0$ and $a \equiv 1 \pmod I$. If in
addition $I \subset \rad(A)$ then $M = 0$.
\end{theorem}

\begin{proof}
Setting $\varphi = 1$ in the previous theorem gives the relation
$a = 1 + a_1 + \cdots + a_n = 0$ as endomorphisms of $M$, that is $aM = 0$, and
$a \equiv 1 \pmod I$. If $I \subset \rad(A)$ then $a$ is a unit of $A$, so that on multiplying
both sides of $aM = 0$ by $a^{-1}$ we get $M = 0$.
\end{proof}

\begin{remark}
This theorem is usually referred to as Nakayama's lemma, but
the late Professor Nakayama maintained that it should be referred to as
a theorem of Krull and Azumaya; it is in fact difficult to determine which
of these three first had the result in the case of commutative rings, so we
refer to it as NAK in this book. Of course, this result can easily be proved
without using determinants, by induction on the number of generators
of $M$.
\end{remark}

\begin{corollary*}
Let $A$ be a ring and $I$ an ideal contained in $\rad(A)$. Suppose
that $M$ is an $A$-module and $N \subset M$ a submodule such that $M/N$ is finite
over $A$. Then $M = N + IM$ implies $M = N$.
\end{corollary*}

\begin{proof}
Setting $\bar{M} = M/N$ we have $\bar{M} = I\bar{M}$, so that, by the
theorem, $\bar{M} = 0$.
\end{proof}

If $W$ is a set of generators of an $A$-module $M$ which is minimal, in the
sense that any proper subset of $W$ does not generate $M$, then $W$ is said
to be a \emph{minimal basis}\index{minimal basis} of $M$. Two minimal bases do not necessarily have
the same number of elements; for example, when $M = A$, if $x$ and $y$ are
non-units of $A$ such that $x + y = 1$ then both $\{1\}$ and $\{x, y\}$ are
minimal bases of $A$. However, if $A$ is a local ring then the situation is clear:

\begin{theorem}\label{thm:2.3}
Let $(A, m, k)$ be a local ring and $M$ a finite $A$-module; set
$\bar{M} = M/mM$. Now $\bar{M}$ is a finite-dimensional vector space over $k$, and we
write $n$ for its dimension. Then:
\begin{enumerate}[(i)]
  \item If we take a basis $\{\bar{\omega}_1, \ldots, \bar{\omega}_n\}$ for $\bar{M}$ over $k$, and choose an
        inverse image $\omega_i \in M$ of each $\bar{\omega}_i$ then $\{\omega_1, \ldots, \omega_n\}$ is a minimal basis
        of $M$;
  \item conversely every minimal basis of $M$ is obtained in this way, and
        so has $n$ elements.
  \item If $\{\omega_1, \ldots, \omega_n\}$ and $\{\omega_1', \ldots, \omega_n'\}$ are both minimal bases of $M$, and
        $\omega_i = \sum_j a_{ij}\omega_j'$ with $a_{ij} \in A$ then $\det(a_{ij})$ is a unit of $A$, so that $(a_{ij})$ is an
        invertible matrix.
\end{enumerate}
\end{theorem}

\begin{proof}
(i) $M = \sum A\omega_i + mM$, and $M$ is finitely generated (hence also
$M/\sum A\omega_i$), so that by the above corollary $M = \sum A\omega_i$. If $\{\omega_1, \ldots, \omega_n\}$ is
not minimal, so that, for example, $\{\omega_2, \ldots, \omega_n\}$ already generates $M$ then
$\{\bar{\omega}_2, \ldots, \bar{\omega}_n\}$ generates $\bar{M}$, which is a contradiction. Hence $\{\omega_1, \ldots, \omega_n\}$
is a minimal basis.

(ii) If $\{\omega_1, \ldots, \omega_n\}$ is a minimal basis of $M$ and we set $\bar{\omega}_i$ for the image
of $\omega_i$ in $\bar{M}$, then $\bar{\omega}_1, \ldots, \bar{\omega}_n$ generate $\bar{M}$, and are linearly independent
over $k$; indeed, otherwise some proper subset of $\{\bar{\omega}_1, \ldots, \bar{\omega}_n\}$ would be
a basis of $\bar{M}$, and then by (i) a proper subset of $\{\omega_1, \ldots, \omega_n\}$ would
generate $M$, which is a contradiction.

(iii) Write $\bar{a}_{ij}$ for the image in $k$ of $a_{ij}$, so that $\bar{\omega}_i = \sum \bar{a}_{ij}\bar{\omega}_j$ holds
in $\bar{M}$. Since $(\bar{a}_{ij})$ is the matrix transforming one basis of the vector space
$\bar{M}$ into another, its determinant is non-zero. Since $\det(a_{ij}) \bmod m =
\det(\bar{a}_{ij}) \neq 0$ it follows that $\det(a_{ij})$ is a unit of $A$. By Cramer's formula
the inverse matrix of $(a_{ij})$ exists as a matrix with entries in $A$.
\end{proof}

We give another interesting application of NAK, the proof of which is
due to Vasconcelos~\cite{Vas2}.

\begin{theorem}\label{thm:2.4}
Let $A$ be a ring and $M$ a finite $A$-module. If $f\colon M \to M$
is an $A$-linear map and $f$ is surjective then $f$ is also injective, and is thus
an automorphism of $M$.
\end{theorem}

\begin{proof}
Since $f$ commutes with scalar multiplication by elements of $A$, we
can view $M$ as an $A[X]$-module by setting $X \cdot m = f(m)$ for $m \in M$. Then by
assumption $XM = M$, so that by NAK there exists $Y \in A[X]$ such that
$(1 + XY)M = 0$. Now for $u \in \Ker(f)$ we have $0 = (1 + XY)(u) = u + Yf(u) = u$, so
that $f$ is injective.
\end{proof}

\begin{theorem}\label{thm:2.5}
Let $(A, m)$ be a local ring; then a projective module over $A$
is free (for the definition of projective module, see Appendix B, p.~277).
\end{theorem}

\begin{proof}
\renewcommand{\qedsymbol}{}% 原书此证明延续至 Lemma 1/2 及文末, 中间无 QED
This is easy when $M$ is finite: choose a minimal basis $\omega_1, \ldots, \omega_n$ of $M$
and define a surjective map $\varphi\colon F \to M$ from the free module
$F = Ae_1 \oplus \cdots \oplus Ae_n$ to $M$ by $\varphi(\sum a_ie_i) = \sum a_i\omega_i$. If we set $K = \Ker(\varphi)$ then, from
the minimal basis property,
\[
  \sum a_i\omega_i = 0 \Rightarrow a_i \in m \text{ for all } i.
\]
Thus $K \subset mF$. Because $M$ is projective, there exists $\psi\colon M \to F$ such that
$F = \psi(M) \oplus K$, and it follows that $K = mK$. On the other hand, $K$ is a
quotient of $F$, therefore finite over $A$, so that $K = 0$ by NAK and $F \cong M$.

The result was proved by Kaplansky~\cite{Kap2} without the assumption that
$M$ is finite. He proves first of all the following lemma, which holds for
any ring (possibly non-commutative).
\end{proof}

\begin{lemma}\label{lem:2.1}
Let $R$ be any ring, and $F$ an $R$-module which is a direct sum of
countably generated submodules; if $M$ is an arbitrary direct summand of $F$
then $M$ is also a direct sum of countably generated submodules.
\end{lemma}

\begin{proof}[Proof of Lemma 1]
\renewcommand{\qedsymbol}{}% 同上: 原书无中间 QED
Suppose that $F = M \oplus N$, and that $F = \bigoplus_{\lambda \in \Lambda} E_\lambda$, where
each $E_\lambda$ is countably generated. By transfinite induction, we construct a
well-ordered family $\{F_\alpha\}$ of submodules of $F$ with the following properties:
\begin{enumerate}[(i)]
  \item if $\alpha < \beta$ then $F_\alpha \subset F_\beta$,
  \item $F = \bigcup_\alpha F_\alpha$,
  \item if $\alpha$ is a limiting ordinal then $F_\alpha = \bigcup_{\beta < \alpha} F_\beta$,
  \item $F_{\alpha+1}/F_\alpha$ is countably generated,
  \item $F_\alpha = M_\alpha \oplus N_\alpha$, where $M_\alpha = M \cap F_\alpha$, $N_\alpha = N \cap F_\alpha$,
  \item each $F_\alpha$ is a direct sum of $E_\lambda$ taken over a suitable subset of $\Lambda$.
\end{enumerate}
We now construct such a family $\{F_\alpha\}$. Firstly, set $F_0 = (0)$. For an
ordinal $\alpha$, assume that $F_\beta$ has been defined for all ordinals $\beta < \alpha$. If $\alpha$ is
a limiting ordinal, set $F_\alpha = \bigcup_{\beta<\alpha} F_\beta$. If $\alpha$ is of the form $\alpha = \beta + 1$, let $Q_1$
be any one of the $E_\lambda$ not contained in $F_\beta$ (if $F_\beta = F$ then the construction
stops at $F_\beta$). Take a set $x_{11}, x_{12}, \ldots$ of generators of $Q_1$, and decompose $x_{11}$
into its $M$- and $N$-components; now let $Q_2$ be the direct sum of the
finitely many $E_\lambda$ which are necessary to write each of these two components
in the decomposition $F = \bigoplus E_\lambda$, and let $x_{21}, x_{22}, \ldots$ be generators of $Q_2$.
Next decompose $x_{12}$ into its $M$- and $N$-components, let $Q_3$ be the direct
sum of the finitely many $E_\lambda$ needed to write these components, and let
$x_{31}, x_{32}, \ldots$ be generators of $Q_3$. Then carry out the same procedure with
$x_{21}$, getting $x_{41}, x_{42}, \ldots$, then do the same for $x_{13}$. Carrying out the same
procedure for each of the $x_{ij}$ in the order $x_{11}, x_{12}, x_{21}, x_{13}, x_{22}, x_{31}, \ldots$ we get
countably many elements $x_{ij}$. We let $F_\alpha$ be the submodule of $F$ generated
by $F_\beta$ and the $x_{ij}$, and this satisfies all our requirements. This gives the
family $\{F_\alpha\}$.

Now $M = \bigcup_\alpha M_\alpha$, with each $M_\alpha$ a direct summand of $F_\alpha$, and
$M_{\alpha+1} \supset M_\alpha$, so that $M_\alpha$ is also a direct summand of $M_{\alpha+1}$. Moreover,
\[
  F_{\alpha+1}/F_\alpha = (M_{\alpha+1}/M_\alpha) \oplus (N_{\alpha+1}/N_\alpha),
\]
and hence $M_{\alpha+1}/M_\alpha$ is countably generated. Thus we can write
\[
  M_{\alpha+1} = M_\alpha \oplus M'_{\alpha+1}, \quad \text{with } M'_{\alpha+1} \text{ countably generated}.
\]
When $\alpha$ is a limit ordinal, since $M_\alpha = \bigcup_{\beta<\alpha} M_\beta$, we set $M'_\alpha = 0$. Then finally
we can write
\[
  M = \bigoplus M'_\alpha, \quad \text{with } M'_\alpha \text{ (at most) countably generated.}
\]
\end{proof}

Of course a free module satisfies the assumption of Lemma~\hyperref[lem:2.1]{1}, so that,
in particular, we see that any projective module is a direct sum of countably
generated projective modules. Thus in the proof of Theorem~\ref{thm:2.5} we can
assume that $M$ is countably generated.

\begin{lemma}\label{lem:2.2}
Let $M$ be a projective module over a local ring $A$, and $x \in M$.
Then there exists a direct summand of $M$ containing $x$ which is a free
module.
\end{lemma}

\begin{proof}[Proof of Lemma 2]
\renewcommand{\qedsymbol}{}% 同上: 原书无中间 QED
We write $M$ as a direct summand of a free module
$F = M \oplus N$. Choose a basis $B = \{u_i\}_{i \in I}$ of $F$ such that the given element
$x$ has the minimum possible number of non-zero coordinates when
expressed in this basis. Then if $x = u_1a_1 + \cdots + u_na_n$ with $0 \neq a_i \in A$, we have
\[
  a_i \notin \sum_{j \neq i} Aa_j \quad \text{for } i = 1, 2, \ldots, n;
\]
indeed, if, say, $a_1 = \sum_{i=2}^n b_ia_i$ then $x = \sum_{i=2}^n (u_i + u_1b_i)a_i$, which contradicts
the choice of $B$. Now set $u_i = y_i + z_i$ with $y_i \in M$ and $z_i \in N$; then
\[
  x = \sum a_iu_i = \sum a_iy_i.
\]
If we write $y_i = \sum_{j=1}^n c_{ij}u_j + t_i$, with $t_i$ linear combinations of elements of
$B$ other than $u_1, \ldots, u_n$, we get relations $a_i = \sum_{j=1}^n a_jc_{ji}$, and, hence, in view
of what we have seen above, we must have
\[
  1 - c_{ii} \in m \quad \text{and} \quad c_{ij} \in m \text{ for } i \neq j.
\]
It follows that the matrix $(c_{ij})$ has an inverse (this can be seen from the
fact that the determinant is $\equiv 1 \pmod m$, or by elimination). Thus replacing
$u_1, \ldots, u_n$ by $y_1, \ldots, y_n$ in $B$, we still have a basis of $F$. Hence, $F_1 = \sum y_iA$ is
a direct summand of $F$, and hence also of $M$, and satisfies all the
requirements of Lemma~\hyperref[lem:2.2]{2}.
\end{proof}

To prove the theorem, let $M$ be a countably generated projective module,
$M = A\omega_1 + A\omega_2 + \cdots$. By Lemma~\hyperref[lem:2.2]{2}, there exists a free module $F_1$ such that
$\omega_1 \in F_1$, and $M = F_1 \oplus M_1$, where $M_1$ is a projective module. Let $\omega_2'$
be the $M_1$-component of $\omega_2$ in the decomposition $M = F_1 \oplus M_1$, and take
a free module $F_2$ such that $\omega_2' \in F_2$ and $M_1 = F_2 \oplus M_2$, where $M_2$ is a
projective module. Let $\omega_3'$ be the $M_2$-component of $\omega_3$ in $M = F_1 \oplus F_2 \oplus M_2$;
proceeding in the same way, we get
\[
  M = F_1 \oplus F_2 \oplus \cdots,
\]
so that $M$ is a free module. $\blacksquare$

We say that an $A$-module $M \neq 0$ is a \emph{simple module}\index{simple module} if it has no
submodules other than $0$ and $M$ itself. For any $0 \neq \omega \in M$, we then have
$M = A\omega$. Now $A\omega \cong A/\ann(\omega)$, but in order for this to be simple, $\ann(\omega)$
must be a maximal ideal of $A$. Hence, any simple $A$-module is isomorphic
to $A/m$ with $m$ a maximal ideal, and conversely an $A$-module of this form
is simple. If $M$ is an $A$-module, a chain
\[
  M = M_0 \supset M_1 \supset \cdots \supset M_r = 0
\]
of submodules of $M$ is called a \emph{composition series}\index{composition series} of $M$ if every $M_i/M_{i+1}$ is
simple; $r$ is called the \emph{length} of the composition series. If a composition series
of $M$ exists, its length is an invariant of $M$ independent of the choice of
composition series. More precisely, if $M$ has a composition series of length
$r$, and if $M_1 \supset N_1 \supset \cdots \supset N_s$ is a strictly descending chain of sub-
modules, then we have $s < r$. This invariance corresponds to part of the
basic Jordan--H\"older theorem in group theory, but it can easily be proved
on its own by induction, and the reader might like to do this as an exercise.
The length of a composition series of $M$ is called the \emph{length of $M$}\index{length of a module}, and written
$l(M)$; if $M$ does not have a composition series we set $l(M) = \infty$. A necessary
and sufficient condition for the existence of a composition series of $M$ is that
the submodules of $M$ should satisfy both the ascending and descending
chain conditions (for which see \S\ref{sec:3}). In general, if $N \subset M$ is a submodule,
we have
\[
  l(M) = l(N) + l(M/N).
\]
If $0 \to M_1 \to M_2 \to \cdots \to M_r \to 0$ is an exact sequence of $A$-modules and
each $M_i$ has finite length then
\[
  \sum_{i=1}^{r} (-1)^i l(M_i) = 0.
\]

If $m$ is a maximal ideal of $A$ and is finitely generated over $A$ then
$l(A/m^\nu) < \infty$. In fact,
\[
  l(A/m^\nu) = l(A/m) + l(m/m^2) + \cdots + l(m^{\nu-1}/m^\nu);
\]
now each $m^i/m^{i+1}$ is a finite-dimensional vector space over the field
$k = A/m$, and since its $A$-submodules are the same thing as its vector
subspaces, $l(m^i/m^{i+1})$ is equal to the dimension of $m^i/m^{i+1}$ as $k$-vector
space. (This shows that $A/m^\nu$ is an Artinian ring, see \S\ref{sec:3}.)

Considering $l(A/m^\nu)$ for all $\nu$, we get a function of $\nu$ which is intimately
related to the ring structure of $A$, and which also plays a role in the resolution
of singularities in algebraic and complex analytic geometry; this is studied in
Chapter~\ref{ch:5}.

We say that an $A$-module $M$ is \emph{of finite presentation}\index{module of finite presentation} if there exists an
exact sequence of the form
\[
  A^p \to A^q \to M \to 0.
\]
This means that $M$ can be generated by $q$ elements $\omega_1, \ldots, \omega_q$ in such a way
that the module $R = \{(a_1, \ldots, a_q) \in A^q \mid \sum a_i\omega_i = 0\}$ of linear relations
holding between the $\omega_i$ can be generated by $p$ elements.

\begin{theorem}\label{thm:2.6}
Let $A$ be a ring, and suppose that $M$ is an $A$-module of finite
presentation. If
\[
  0 \to K \to N \to M \to 0
\]
is an exact sequence and $N$ is finitely generated then so is $K$.
\end{theorem}

\begin{proof}
By assumption there exists an exact sequence of the form
$L_2 \xrightarrow{\alpha} L_1 \xrightarrow{\varphi} M \to 0$, where $L_1$ and $L_2$ are free modules of finite rank.
From this we get the following commutative diagram (see Appendix B):
\[
  \begin{tikzcd}[column sep=large, row sep=large]
    & L_2 \arrow[r, "\alpha"] \arrow[d, "\beta"] & L_1 \arrow[r, "\varphi"] \arrow[d, "u"] & M \arrow[r] \arrow[d, equal] & 0 \\
    0 \arrow[r] & K \arrow[r, "\gamma"] & N \arrow[r, "f"] & M \arrow[r] & 0
  \end{tikzcd}
\]
If we write $N = At_1 + \cdots + At_q$, then there exist $u_i \in L_1$ such that
$\varphi(u_i) = f(t_i)$. Set $t_i' = t_i - u(u_i)$; then $f(t_i') = 0$, so that we can write
$t_i' = \gamma(\eta_i)$ with $\eta_i \in K$. Let us now prove that
\[
  K = \beta(L_2) + A\eta_1 + \cdots + A\eta_q.
\]
For any $\eta \in K$, set $\psi(\eta) = \sum a_i t_i$; then
\[
  \psi\left(\eta - \sum a_i t_i'\right) = \sum a_i (t_i - t_i') = u\left(\sum a_i u_i\right),
\]
and since $0 = f\left(\sum a_i t_i\right) = \varphi\left(\sum a_i u_i\right)$, we can write $\sum a_i u_i = \alpha(w)$
with $w \in L_2$. Now
\[
  \eta - \sum a_i t_i' = u\left(\sum a_i u_i\right) = \beta(w),
\]
so that $\eta = \beta(w) + \sum a_i\eta_i$, and this proves our assertion.
\end{proof}

% ============ 习题 ============
\section*{Exercises to \S\ref{sec:2}}
\emph{Prove the following propositions.}

\begin{exercise}\label{ex:2.1-ex}
Let $A$ be a ring and $I$ a finitely generated ideal satisfying $I = I^2$; then $I$ is
generated by an idempotent $e$ (an element $e$ satisfying $e^2 = e$).
\end{exercise}

\begin{exercise}\label{ex:2.2-ex}
Let $A$ be a ring, $I$ an ideal of $A$ and $M$ a finite $A$-module; then
\[
  \sqrt{\ann(M/IM)} = \sqrt{\ann(M) + I}.
\]
\end{exercise}

\begin{exercise}\label{ex:2.3-ex}
Let $M$ and $N$ be submodules of an $A$-module $L$. If $M + N$ and $M \cap N$ are
finitely generated then so are $M$ and $N$.
\end{exercise}

\begin{exercise}\label{ex:2.4-ex}
Let $A$ be a (commutative) ring, $A \neq 0$. An $A$-module is said to be \emph{free of rank $n$}
if it is isomorphic to $A^n$.
\begin{enumerate}[(a)]
  \item If $A^n \cong A^m$ then $n = m$; prove this by reducing to the case of a field.
        (Note that there are counter-examples to this for non-commutative rings.)
  \item Let $C = (c_{ij})$ be an $n \times m$ matrix over $A$, and suppose that $C$ has a non-
        zero $r \times r$ minor, but that all the $(r + 1) \times (r + 1)$ minors are $0$. Show then
        that if $r < m$, the $m$ column vectors of $C$ are linearly dependent. (Hint: you
        can assume that $m = r + 1$.) Deduce from this an alternative proof of (a).
  \item If $A$ is a local ring, any minimal basis of the free module $A^n$ is a basis
        (that is, a linearly independent set of generators).
\end{enumerate}
\end{exercise}

\begin{exercise}\label{ex:2.5-ex}
Let $A$ be a ring, and $0 \to L \to M \to N \to 0$ an exact sequence of $A$-
modules.
\begin{enumerate}[(a)]
  \item If $L$ and $N$ are both of finite presentation then so is $M$.
  \item If $L$ is finitely generated and $M$ is of finite presentation then $N$ is of
        finite presentation.
\end{enumerate}
\end{exercise}
